\begin{figure}[htbp]
  \centering
  \includegraphics[width=160mm]{q2_posterior_general.pdf}
  \caption{非圆心检测点处，不同首次示向度对应的后验位置密度}
  \label{fig:q2-posterior}
\end{figure}

\begin{figure}[htbp]
  \centering
  \includegraphics[width=130mm]{q2_feedback_regions.pdf}
  \caption{两种示向度反馈下的剩余区域与完成条件}
  \label{fig:q2-feedback-regions}
\end{figure}

\begin{figure}[htbp]
  \centering
  \includegraphics[width=150mm]{q2_feedback_distribution.pdf}
  \caption{第二次示向度的预测密度及满足完成条件的反馈}
  \label{fig:q2-feedback}
\end{figure}

\begin{figure}[htbp]
  \centering
  \includegraphics[width=160mm]{q2_candidates.pdf}
  \caption{11个权重下的候选检测点及评价指标（\(a=0,\beta=0\)）。颜色表示完成概率权重，虚线连接相邻权重的数值结果。}
  \label{fig:q2-candidates}
\end{figure}